\section{Introducción: por qué estos informes apuntan a los Agents}

Esta sección establece primero los objetivos de aprendizaje de todo el episodio. El programa revisa varios de los materiales técnicos sobre Agents más dignos de una lectura atenta de los últimos meses: el informe técnico de Kimi K2, el anuncio de lanzamiento de ChatGPT Agent, la entrada técnica del blog de Qwen3-Coder y la experiencia de Manus en ingeniería de contexto. Aunque parecen venir de empresas y productos distintos, todos apuntan a un mismo tema: los modelos grandes están pasando de «responder preguntas» a «completar tareas dentro de un entorno».

Lo esencial de un Agent no es solo llamar herramientas, sino el ciclo completo: percibir el entorno, decidir, actuar, leer la retroalimentación y corregir la trayectoria. Este ciclo trae consigo todo un conjunto de problemas de ingeniería: datos sintéticos, generación de trayectorias, aprendizaje por refuerzo, recompensas verificables, entornos aislados de seguridad, ingeniería de contexto, KV cache, selección de herramientas, memoria en el sistema de archivos y organización de múltiples Agents.

\begin{figure}[H]
\centering
\includegraphics[width=0.96\textwidth]{report-comparison.png}
\caption{Comparación de los cuatro informes: Kimi, ChatGPT Agent, Qwen3-Coder y Manus enfatizan cada uno una capa distinta. Diagrama conceptual de elaboración propia, organizado a partir de la estructura del video y de los materiales oficiales.}
\end{figure}

\begin{knowledgebox}{Lectura de la figura: cada material aborda una capa}
Kimi K2 aborda el modelo abierto y el agentic training; ChatGPT Agent, el sistema de producto unificado; Qwen3-Coder, el RL y el uso de herramientas en entornos de código; Manus, la ingeniería de contexto de los Agents en producción. En conjunto forman la cadena completa de los Agents: del modelo y el entrenamiento al producto y el despliegue de ingeniería.
\end{knowledgebox}

\begin{importantbox}{Tesis central del episodio}
El progreso de los Agents no consiste en mejorar la capacidad de un solo modelo, sino en ingeniería de sistemas: solo optimizando en conjunto el modelo, el entorno, las tareas, las herramientas, las recompensas, la seguridad y el contexto es posible llevar la inteligencia de producir texto a completar tareas reales.
\end{importantbox}

\begin{knowledgebox}{Orden de estudio de los cuatro materiales}
\begin{tabularx}{\textwidth}{p{0.20\textwidth}p{0.34\textwidth}X}
\toprule
Material & Qué revisar primero & Valor didáctico \\
\midrule
Kimi K2 & MoE, MuonClip, agentic data, joint RL & Ver cómo se entrena un modelo abierto para volverlo Agentic. \\
ChatGPT Agent & Operator + Deep Research + computadora con herramientas & Ver cómo un Agent se convierte en un producto para usuarios. \\
Qwen3-Coder & Code RL, long-horizon RL, Qwen Code & Ver por qué los entornos de código son adecuados para el RL de Agents. \\
Manus & KV cache, enmascaramiento de herramientas, sistema de archivos, conservación de errores & Ver cómo un Agent en producción logra funcionar de forma estable gracias a la ingeniería de contexto. \\
\bottomrule
\end{tabularx}
\end{knowledgebox}

\begin{knowledgebox}{Resumen de cifras clave}
\begin{tabularx}{\textwidth}{p{0.22\textwidth}p{0.34\textwidth}X}
\toprule
Material & Cifras o hechos clave & Implicación didáctica \\
\midrule
Kimi K2 & Cerca de 1T de parámetros totales, 32B de parámetros activos, 15.5T tokens & Los modelos MoE abiertos también entraron al entrenamiento de sistemas a escala muy grande. \\
Kimi K2 & SWE-Bench Verified 65.8, Tau2-Bench 66.1 & Los Agentic benchmarks se vuelven indicadores centrales de los informes de modelos. \\
Qwen3-Coder & 480B MoE, 35B active, 256K/1M context & Un Agent de código necesita un contexto muy largo y comprensión a nivel de repositorio. \\
Qwen3-Coder & 20,000 entornos paralelos para long-horizon RL & El cuello de botella del RL de Agents es escalar los entornos. \\
Manus & Proporción de tokens de entrada/salida de aproximadamente 100:1 & La KV cache es crucial para el costo y la latencia. \\
\bottomrule
\end{tabularx}
\end{knowledgebox}

\subsection{Resumen de la sección}

EP110 es la clase sobre informes técnicos de Agents. Conecta la teoría de Agents de EP115, la perspectiva de empresa de modelos sobre K2 de EP113 y el Agentic Model de nivel empresarial de EP116, y muestra «el poder de la ingeniería de sistemas».
